\chapter{Simple Data}

This chapter introduces the atomic (non-compound) data types beyond numbers ---
booleans, symbols, strings, and characters --- together with the conditional
expression \code{cond}, which lets a function choose among cases.

\section{Booleans and predicates}

The two Boolean values are \code{true} and \code{false} (also written
\code{\#true}/\code{\#t} and \code{\#false}/\code{\#f}). A \textbf{predicate}
is a function that produces a Boolean; by convention its name ends in \code{?}.

\begin{itemize}
  \item Comparisons: \code{=}, \code{<}, \code{<=}, \code{>}, \code{>=} on numbers.
  \item Type predicates: \code{number?}, \code{integer?}, \code{boolean?}, \code{symbol?}, \code{string?}, \code{char?}.
  \item Equality: \code{equal?} works on any two values; \code{symbol=?}, \code{string=?}, \code{char=?} are type-specific.
\end{itemize}

\subsection{\code{and}, \code{or}, \code{not}}
\code{not} is an ordinary function. \code{and} and \code{or} are
\textbf{special forms}: they evaluate their arguments left to right and stop as
soon as the answer is known (\textbf{short-circuit evaluation}).

\begin{important}[Substitution rules for and/or]
\begin{itemize}
  \item \code{(and false ...)} $\Rightarrow$ \code{false}
  \item \code{(and true e2 ...)} $\Rightarrow$ \code{(and e2 ...)}
  \item \code{(and)} $\Rightarrow$ \code{true}
  \item \code{(or true ...)} $\Rightarrow$ \code{true}
  \item \code{(or false e2 ...)} $\Rightarrow$ \code{(or e2 ...)}
  \item \code{(or)} $\Rightarrow$ \code{false}
\end{itemize}
\end{important}

Short-circuiting lets us guard a dangerous test:
\begin{lstlisting}[style=racket]
(and (not (= x 0)) (> (/ 1 x) 2))   ; never divides by zero
\end{lstlisting}

\section{Conditional expressions}

\begin{lstlisting}[style=racket]
(cond
  [question1 answer1]
  [question2 answer2]
  ...
  [else answerk])
\end{lstlisting}
Each \code{[question answer]} pair is a \textbf{clause}. Questions are tried in
order; the answer of the first question that evaluates to \code{true} is the
result.

\begin{important}[Substitution rules for cond]
\begin{itemize}
  \item \code{(cond [false e] ...)} $\Rightarrow$ \code{(cond ...)}
  \item \code{(cond [true e] ...)} $\Rightarrow$ \code{e}
  \item \code{(cond [else e])} $\Rightarrow$ \code{e}
\end{itemize}
Only the \emph{first} question is evaluated at each step; answers of
unchosen clauses are never evaluated. If no question is true and there is no
\code{else}, it is a run-time error.
\end{important}

\begin{example}[Grading]
\begin{lstlisting}[style=racket]
;; (letter-grade mark) produces the letter grade for a percentage mark.
;; Examples:
(check-expect (letter-grade 95) 'A)
(check-expect (letter-grade 42) 'F)

;; letter-grade: Num -> Sym
;;   requires: 0 <= mark <= 100
(define (letter-grade mark)
  (cond [(>= mark 80) 'A]
        [(>= mark 70) 'B]
        [(>= mark 60) 'C]
        [(>= mark 50) 'D]
        [else 'F]))

;; Tests (boundaries of every interval):
(check-expect (letter-grade 80) 'A)
(check-expect (letter-grade 79) 'B)
(check-expect (letter-grade 70) 'B)
(check-expect (letter-grade 60) 'C)
(check-expect (letter-grade 50) 'D)
(check-expect (letter-grade 49.9) 'F)
\end{lstlisting}
Because clauses are tried in order, the second clause need not say
\code{(and (< mark 80) (>= mark 70))}. Take advantage of this, but be careful
when reordering.
\end{example}

The questions split the range of marks into intervals. Each boundary value
belongs to the higher grade, and those boundaries are exactly what the tests
check:

\begin{center}
\begin{tikzpicture}[x=1.2mm,y=1mm]
  \foreach \a/\b/\g/\c in {0/50/F/red!25,50/60/D/orange!30,60/70/C/yellow!40,70/80/B/green!25,80/100/A/blue!20} {
    \fill[\c] (\a,0) rectangle (\b,6);
    \node at ({(\a+\b)/2},3) {\code{'\g}};
  }
  \draw (0,0) rectangle (100,6);
  \foreach \t in {0,50,60,70,80,100} {
    \draw (\t,0) -- (\t,-1.5) node[below,font=\scriptsize] {\t};
  }
  \foreach \t in {50,60,70,80} \fill (\t,0) circle (0.6mm);
  \draw[-{Stealth},thick,blue!50!black] (75,14) node[above,font=\scriptsize]{\code{(letter-grade 75)}} -- (75,7);
\end{tikzpicture}
\end{center}

Read as a flowchart, the \code{cond} asks its questions one at a time and
stops at the first \code{true}:

\begin{center}
\begin{tikzpicture}[q/.style={draw,diamond,aspect=1.6,inner sep=1pt,fill=blue!6,font=\scriptsize},
  a/.style={draw,rounded corners,font=\ttfamily\small,minimum width=8mm,minimum height=6mm},
  arr/.style={-{Stealth},thick}, lbl/.style={font=\scriptsize}]
  \node[q] (q1) at (0,0) {$\ge 80$?};
  \node[q] (q2) at (2.6,0) {$\ge 70$?};
  \node[q] (q3) at (5.2,0) {$\ge 60$?};
  \node[q] (q4) at (7.8,0) {$\ge 50$?};
  \node[a] (aF) at (10.3,0) {'F};
  \node[a] (aA) at (0,-1.6) {'A};
  \node[a] (aB) at (2.6,-1.6) {'B};
  \node[a] (aC) at (5.2,-1.6) {'C};
  \node[a] (aD) at (7.8,-1.6) {'D};
  \node[lbl,left=3mm of q1] (m) {\code{mark}};
  \draw[arr] (m) -- (q1);
  \foreach \x/\y in {q1/q2,q2/q3,q3/q4} \draw[arr] (\x) -- node[lbl,above]{no} (\y);
  \draw[arr] (q4) -- node[lbl,above]{else} (aF);
  \foreach \x/\y in {q1/aA,q2/aB,q3/aC,q4/aD} \draw[arr] (\x) -- node[lbl,right]{yes} (\y);
\end{tikzpicture}
\end{center}

\begin{example}[Trace of cond]
\begin{lstlisting}[style=racket]
(letter-grade 75)
=> (cond [(>= 75 80) 'A] [(>= 75 70) 'B] [(>= 75 60) 'C] [(>= 75 50) 'D] [else 'F])
=> (cond [false 'A] [(>= 75 70) 'B] [(>= 75 60) 'C] [(>= 75 50) 'D] [else 'F])
=> (cond [(>= 75 70) 'B] [(>= 75 60) 'C] [(>= 75 50) 'D] [else 'F])
=> (cond [true 'B] [(>= 75 60) 'C] [(>= 75 50) 'D] [else 'F])
=> 'B
\end{lstlisting}
\end{example}

\subsection{Simplifying conditionals}
\begin{itemize}
  \item \code{(cond [test true] [else false])} is just \code{test}.
  \item Nested \code{cond}s whose clauses can be merged with \code{and} often read better flattened.
  \item A function that produces a Boolean rarely needs \code{cond} at all; combine \code{and}/\code{or}/\code{not}.
\end{itemize}

\subsection{Testing conditionals}
At least one test per clause, plus tests at each boundary between intervals.
For a \code{cond} with $k$ clauses expect at least $k$ tests; boundary values
are where off-by-one bugs live.

\section{Symbols}
A \textbf{symbol} is written with a leading quote: \code{'blue},
\code{'CS135}. Symbols are atomic --- they cannot be taken apart --- and the
only useful operation is comparison with \code{symbol=?} (or \code{equal?}).
They are ideal for representing a small fixed set of options, like an
enumeration:

\begin{lstlisting}[style=racket]
;; A TrafficLight is (anyof 'red 'yellow 'green)

;; (next-light c) produces the light shown after c.
;; next-light: TrafficLight -> TrafficLight
(check-expect (next-light 'red) 'green)
(define (next-light c)
  (cond [(symbol=? c 'red) 'green]
        [(symbol=? c 'green) 'yellow]
        [(symbol=? c 'yellow) 'red]))
\end{lstlisting}

\section{Strings}
A \textbf{string} is a sequence of characters in double quotes: \code{"CS 135"}.
Unlike symbols, strings have structure and many operations:
\begin{center}
\begin{tabular}{ll}
\toprule
Function & Example \\
\midrule
\code{string-length} & \code{(string-length "abc")} $\Rightarrow$ 3 \\
\code{string-append} & \code{(string-append "ab" "cd")} $\Rightarrow$ \code{"abcd"} \\
\code{substring} & \code{(substring "hello" 1 3)} $\Rightarrow$ \code{"el"} \\
\code{string=?}, \code{string<?} & lexicographic comparison \\
\code{number->string} & \code{(number->string 42)} $\Rightarrow$ \code{"42"} \\
\code{string->list} & \code{(string->list "ab")} $\Rightarrow$ \code{(list \#\textbackslash a \#\textbackslash b)} \\
\code{list->string} & inverse of the above \\
\bottomrule
\end{tabular}
\end{center}
Note that \code{substring} uses a half-open range: start index included, end
index excluded, with indices starting at 0.

\paragraph{Symbol or string?} Use a symbol when the value is a label you only
compare (colours, suits, states). Use a string when the value is text that
might be examined, combined, or displayed.

\section{Characters}
A \textbf{character} is written \code{\#\textbackslash a}, \code{\#\textbackslash Z},
\code{\#\textbackslash space}, \code{\#\textbackslash newline}. Useful functions:
\begin{itemize}
  \item comparison: \code{char=?}, \code{char<?};
  \item case conversion: \code{char-upcase}, \code{char-downcase};
  \item classification: \code{char-alphabetic?}, \code{char-numeric?}, \code{char-whitespace?};
  \item conversion to and from character codes: \code{char->integer}, \code{integer->char}.
\end{itemize}
Characters become important once strings are converted to lists of characters
(Chapter~\ref{ch:lists}).

\section{Mixed data and type predicates}
Sometimes a function accepts values of several types. The contract uses
\code{anyof}, and the body uses type predicates to distinguish the cases:

\begin{lstlisting}[style=racket]
;; (describe v) produces a string describing v.
;; describe: (anyof Num Sym Str) -> Str
(check-expect (describe 5) "number")
(check-expect (describe 'x) "symbol")
(check-expect (describe "hi") "string")
(define (describe v)
  (cond [(number? v) "number"]
        [(symbol? v) "symbol"]
        [else "string"]))
\end{lstlisting}

\section{Summary}
\begin{itemize}
  \item Booleans, symbols, strings and characters are the atomic types besides numbers.
  \item \code{and}/\code{or} short-circuit; \code{cond} evaluates questions in order and only the chosen answer.
  \item Know the exact substitution rules for \code{cond}, \code{and}, \code{or}.
  \item Test every clause and every boundary.
  \item Symbols for labels, strings for text; \code{anyof} plus type predicates for mixed data.
\end{itemize}
